% =========================================================
% Elsevier journal-ready LaTeX template (elsarticle)
% Paper: "From Models and Algorithms to Decision Makers..."
% =========================================================
% Compile with: pdflatex -> bibtex -> pdflatex -> pdflatex
% (or use Overleaf with elsarticle)
%
% Notes:
% - Explanations are provided as LaTeX comments (%) and won't appear in PDF.
% - Replace journal options (e.g., 1p/3p/5p) depending on target journal.
% - For many Elsevier journals, you may need to use "cas-dc" instead of elsarticle.
%   This template uses elsarticle because it's widely supported and simple.
% =========================================================

\documentclass[preprint,12pt]{elsarticle}
% Options:
%   preprint      : double-spaced, review-friendly
%   review        : (some journals) for review formatting
%   1p / 3p / 5p  : final layout styles (journal-specific)

% ---------------------------------------------------------
% PACKAGES
% ---------------------------------------------------------
\usepackage[T1]{fontenc}
\usepackage[utf8]{inputenc} % If your compiler is pdfLaTeX. For XeLaTeX/LuaLaTeX, this may be unnecessary.
\usepackage{lmodern}

\usepackage{amsmath,amssymb}
\usepackage{graphicx}
\usepackage{booktabs}
\usepackage{multirow}
\usepackage{enumitem}
\usepackage{hyperref}
\usepackage{url}

% ---------------------------------------------------------
% ELSEVIER JOURNAL INFO
% ---------------------------------------------------------
\journal{Journal Name Here}
% % If needed, include Elsevier-specific commands (rare for basic templates).

% ---------------------------------------------------------
% TITLE/AUTHORS
% ---------------------------------------------------------
\begin{document}
	
	\begin{frontmatter}
		
		\title{From Models and Algorithms to Decision Makers: A Systematic Review of Hybrid Augmented Intelligence in the Vehicle Routing Problem}
		
		% % Add affiliations as needed:
		% \author[inst1]{Helen~Surname}
		% \ead{helen.email@domain.edu}
		% \affiliation[inst1]{organization={Universidad Autónoma de Ciudad Juárez},
			%   addressline={...}, city={Ciudad Juárez}, state={Chihuahua}, country={Mexico}}
		
		\author{Author Name(s) Here}
		\affiliation{organization={Institution(s) Here}}
		
		\begin{abstract}
			% ---------------------------------------------------------
			% ABSTRACT GUIDELINES (do not show in PDF)
			% ---------------------------------------------------------
			% Keep it concise (typically 150--250 words).
			% Include:
			% - Context: VRP has moved beyond pure algorithmic optimization toward decision support.
			% - Objective: PRISMA-based systematic review of hybrid augmented intelligence approaches.
			% - Methods: databases, screening criteria, taxonomy construction, time window 2015--2025.
			% - Results: key trends (human-in-the-loop, preference integration, symbolic+subsymbolic AI).
			% - Contribution: taxonomy + synthesis + gaps + future directions.
		\end{abstract}
		
		\begin{keyword}
			% ---------------------------------------------------------
			% KEYWORDS GUIDELINES
			% ---------------------------------------------------------
			% Provide 4--8 keywords.
			Vehicle Routing Problem \sep Augmented Intelligence \sep Human-in-the-loop \sep Hybrid Metaheuristics \sep Preference-based Optimization \sep Decision Support Systems
		\end{keyword}
		
	\end{frontmatter}
	
	% =========================================================
	\section{Introduction}
	% % Purpose:
	% % - Establish the problem: logistics decision-making is not only optimization.
	% % - Motivate hybrid augmented intelligence: humans + algorithms + AI.
	% % - State contributions and research questions.
	
	\subsection{Background: From Optimization Algorithms to Decision-Centered Systems}
	% % Explain evolution:
	% % - Classical OR: exact methods, then heuristics/metaheuristics.
	% % - Modern VRP settings: dynamic constraints, uncertainty, multi-stakeholder goals.
	% % - Shift: from "optimal routes" to "decision-centered systems".
	
	\subsection{Motivation: Why Purely Algorithmic Approaches Are No Longer Enough}
	% % Discuss insufficiency of purely algorithmic VRP solutions:
	% % - Static objectives vs evolving business priorities.
	% % - Multiple criteria and preferences not captured by a single cost function.
	% % - Need for interpretability, trust, and human acceptance.
	% % - Data drift, operational disruptions, and human accountability.
	
	\subsection{Augmented Intelligence and Human-in-the-Loop Optimization}
	% % Define:
	% % - Augmented intelligence: collaboration between AI/optimization and human expertise.
	% % - Human-in-the-loop: interactive preference elicitation, validation, steering.
	% % Mention: interactive optimization, DSS, cognitive embedded models.
	
	\subsection{Objectives and Contributions of This Review}
	% % Clearly state:
	% % - Scope: intersection of VRP + (metaheuristics/optimization) + (MCDM/preferences/HITL) + (symbolic+subsymbolic AI).
	% % - Period: 2015--2025.
	% % - Contributions: (i) PRISMA systematic review, (ii) taxonomy, (iii) synthesis of achievements, (iv) gaps + future directions.
	
	% =========================================================
	\section{Conceptual Foundations}
	% % Purpose:
	% % Align terminology across OR, AI, decision science.
	% % Provide minimal background so readers can follow taxonomy and synthesis.
	
	\subsection{The Vehicle Routing Problem and Its Variants}
	% % Include brief definitions of: VRP, CVRP, VRPTW, MDVRP, stochastic/dynamic VRP, rich VRP.
	% % Explain why VRP is a strong testbed for hybrid/augmented approaches.
	
	\subsection{Metaheuristics and Hybrid Optimization}
	% % Brief overview of metaheuristics used in VRP (GA, ACO, SA, VNS, ALNS, PSO, etc.).
	% % Explain hybridization in OR:
	% % - memetic algorithms, matheuristics, hyper-heuristics
	% % Distinguish from hybrid augmented intelligence (where humans/knowledge are integrated).
	
	\subsection{Symbolic and Subsymbolic AI in Optimization}
	% % Symbolic: rules, constraints, expert systems, knowledge bases, ontologies.
	% % Subsymbolic: ML, deep learning, RL, surrogate models.
	% % Explain why VRP increasingly uses both (e.g., rules for feasibility + ML for prediction).
	
	\subsection{Decision-Making, Preferences, and Human-in-the-Loop Systems}
	% % Introduce:
	% % - MCDM (AHP, TOPSIS, ELECTRE, PROMETHEE) and preference modeling.
	% % - Preference-based optimization / interactive multiobjective optimization.
	% % - DSS and cognitive models.
	% % This section should connect directly to your research questions.
	
	% =========================================================
	\section{Methodology: PRISMA-Based Systematic Review}
	% % Purpose:
	% % Provide full reproducibility: protocol, searches, screening, extraction, synthesis.
	% % Include PRISMA 2020 flow diagram (as figure) in final draft.
	
	\subsection{Review Protocol and Research Questions}
	% % State research questions (RQs) explicitly.
	% % Example set (adapt to your final framing):
	% % RQ1: How are humans + heuristics/metaheuristics + symbolic/subsymbolic AI combined in VRP decision-making?
	% % RQ2: Why are purely algorithmic approaches insufficient in the reviewed settings?
	% % RQ3: Where and how is expert knowledge integrated (rules, preferences, cognitive models)?
	% % RQ4: What outcomes/benefits are reported (solution quality, decision quality, robustness, usability)?
	% % RQ5: What gaps remain (evaluation, benchmarks, reproducibility)?
	
	\subsection{Search Strategy}
	% % Describe databases (e.g., Scopus, Web of Science, IEEE Xplore, ACM, Springer, ScienceDirect).
	% % Provide complete search strings (appendix if long).
	% % Set time window: 2015--2025.
	% % Mention backward/forward snowballing if used.
	
	\subsection{Inclusion and Exclusion Criteria}
	% % Inclusion examples:
	% % - VRP (or variants) AND hybrid approach involving at least two of:
	% %   (metaheuristics/optimization) + (symbolic/subsymbolic AI) + (human interaction/preferences/decision support).
	% % - Peer-reviewed (journals/conferences) depending on scope.
	% % Exclusion examples:
	% % - Purely algorithmic VRP studies with no decision/human/preference/knowledge aspect.
	% % - Non-routing logistics unless explicitly tied to VRP decision support.
	% % Clarify language restrictions (English-only?) if applied.
	
	\subsection{Study Selection Process}
	% % Describe the PRISMA stages:
	% % - Identification -> Screening -> Eligibility -> Included.
	% % Report record counts at each stage.
	% % Mention deduplication method and screening tools (Rayyan, Zotero, etc.) if used.
	
	\subsection{Data Extraction and Classification Framework}
	% % Provide the extraction schema:
	% % - Bibliographic info, VRP variant, objectives (single/multi), uncertainty/dynamics.
	% % - Algorithmic methods (metaheuristic, matheuristic, exact).
	% % - AI components (ML, RL, surrogate, fuzzy, rules, ontology).
	% % - Human role (preference elicitation, steering, validation, tuning).
	% % - Evaluation (benchmarks, real data, user studies, deployment).
	% % Explain how taxonomy categories were derived (deductive, inductive, mixed).
	
	% =========================================================
	\section{Taxonomy of Hybrid Augmented Intelligence Approaches}
	% % Purpose:
	% % Present your main organizing framework and classification of studies.
	% % This should directly answer multiple RQs.
	
	\subsection{Dimensions of Hybridization}
	% % Define your taxonomy axes, e.g.:
	% % (A) Optimization layer: metaheuristics / matheuristics / exact hybrids.
	% % (B) Intelligence layer: symbolic / subsymbolic / fuzzy / preference learning.
	% % (C) Human layer: interactive optimization / DSS / cognitive embedding.
	% % Provide a table summarizing dimensions and codes used in extraction.
	
	\subsection{Role of the Human in the Loop}
	% % Classify human participation modes, e.g.:
	% % - Offline: expert-designed rules, constraint modeling, parameter setting.
	% % - Online: interactive preference elicitation, solution steering, acceptance/rejection feedback.
	% % - Oversight: validation, interpretability, accountability, exception handling.
	
	\subsection{Knowledge Integration Mechanisms}
	% % Explain how expert knowledge enters:
	% % - Rule injection, constraint shaping.
	% % - Fuzzy linguistic preferences and membership functions.
	% % - Learned preference models and surrogate-assisted search.
	% % - Hybrid symbolic–subsymbolic architectures (rules + ML).
	
	\subsection{Architectural Patterns}
	% % Identify patterns commonly found in the literature:
	% % - ML-assisted metaheuristics (predict costs, feasibility, demand, travel times).
	% % - Surrogate models accelerating evaluation.
	% % - Preference-based multiobjective VRP with interactive methods.
	% % - Rule-based controllers orchestrating operators (hyper-heuristics).
	% % - Cognitive/decision-maker models embedded in the loop.
	
	% =========================================================
	\section{Synthesis of Results (2015--2025)}
	% % Purpose:
	% % Summarize findings across studies, not paper-by-paper but theme-by-theme.
	% % Include descriptive statistics: counts by year, method type, VRP variant, evaluation style.
	
	\subsection{Trends in Hybrid AI for VRP}
	% % Discuss temporal evolution:
	% % - Early years: hybrid metaheuristics + fuzzy + DSS.
	% % - Mid period: surrogates, learning-guided search.
	% % - Recent: preference learning, interactive systems, more realistic datasets.
	% % Provide evidence: plots/tables of trends (years vs categories).
	
	\subsection{Reported Achievements}
	% % Summarize what was actually achieved:
	% % - Better trade-off handling (multiobjective, preferences).
	% % - Robustness (uncertainty, dynamics).
	% % - Computational improvements (surrogates, learned heuristics).
	% % - Usability/acceptance (when evaluated).
	% % Distinguish: solution quality vs decision quality.
	
	\subsection{Limitations of Purely Algorithmic Approaches}
	% % Provide synthesized arguments supported by evidence:
	% % - When objectives are uncertain or stakeholder-dependent.
	% % - When interpretability is needed for adoption.
	% % - When real-time changes require human judgment.
	% % Tie back to your motivation and RQ2.
	
	\subsection{Integration of Expert Knowledge}
	% % Identify where expert knowledge appears:
	% % - Model design: constraints, objectives, feasibility rules.
	% % - Search design: operator selection, heuristic rules.
	% % - Decision layer: ranking, preference articulation, acceptance criteria.
	% % Discuss how frequently and how explicitly studies document expert input.
	
	% =========================================================
	\section{Discussion: Toward Decision-Centered Optimization}
	% % Purpose:
	% % Interpret results, connect to broader theory, highlight gaps.
	
	\subsection{From Solving VRPs to Supporting Decisions}
	% % Argue: VRP is increasingly a socio-technical decision problem.
	% % Discuss implications for evaluation, reporting, and system design.
	
	\subsection{Augmented Intelligence vs Full Automation}
	% % Compare the paradigms:
	% % - Automation: maximize objective(s) with minimal human involvement.
	% % - Augmentation: human expertise improves relevance, trust, adaptability.
	% % Discuss explainability, transparency, and accountability.
	
	\subsection{Open Challenges and Research Gaps}
	% % Examples:
	% % - Lack of standardized benchmarks for HITL VRP.
	% % - Weak evaluation of decision quality and human factors (time, cognitive load).
	% % - Reproducibility issues (private datasets, missing code).
	% % - Need for stronger symbolic–subsymbolic integration frameworks.
	% % - Missing cognitive models of decision makers.
	
	% =========================================================
	\section{Future Research Directions}
	% % Provide concrete directions, e.g.:
	% % - Cognitive digital twins for decision makers.
	% % - Adaptive preference learning integrated with metaheuristics.
	% % - Explainable hybrid search and transparent rule+ML fusion.
	% % - Real-time interactive optimization under uncertainty.
	% % - Protocols for evaluating decision quality (beyond objective value).
	
	% =========================================================
	\section{Threats to Validity}
	% % Discuss threats typical for systematic reviews:
	% % - Publication bias and venue bias.
	% % - Terminology fragmentation (papers use different labels for similar ideas).
	% % - Database coverage limitations.
	% % - Screening and coding subjectivity.
	% % - Time window limitations (2015--2025).
	
	% =========================================================
	\section{Conclusions}
	% % Summarize main claims supported by synthesis:
	% % - Hybrid augmented intelligence is increasingly central in VRP.
	% % - Purely algorithmic solutions often fail to capture decision realities.
	% % - Human expertise is integrated via preferences, rules, validation, and interactive loops.
	% % - Important research gaps remain in evaluation and cognitive integration.
	
	% =========================================================
	% REFERENCES
	% =========================================================
	% Elsevier often recommends using natbib with elsarticle.
	% elsarticle loads natbib by default; citations:
	%   \citet{key} for textual, \citep{key} for parenthetical.
	
	\bibliographystyle{elsarticle-num}
	% Alternatives:
	% \bibliographystyle{elsarticle-harv} % author-year
	% \bibliographystyle{elsarticle-num-names}
	
	\bibliography{references}
	
	% =========================================================
	% APPENDICES (optional)
	% =========================================================
	% \appendix
	% \section{Full Search Strings}
	% % Put complete database queries here for reproducibility.
	%
	% \section{Data Extraction Form}
	% % Include your coding schema/table.
	
\end{document}